
\subsubsection{Code Generation Alignment Methods}

SelfCodeAlign achieves 67.1\% pass@1 on HumanEval+ through execution-based feedback (test pass/fail signals). StepCoder similarly uses test-driven training. These methods establish that feedback signal type affects performance but focus exclusively on correctness metrics without measuring cross-dimensional effects.

Preference-based alignment (DPO, RLAIF) trains models on human or AI-generated preference pairs, theoretically capturing holistic code quality. Recent applications to code generation show promise but similarly focus on single-metric benchmarking.

\subsubsection{Multi-Objective Optimization}

PrefGen explicitly trains smart contract generation models for multiple objectives: correctness (pass@k), gas efficiency (gas@k), and security (secure@k). This demonstrates multi-objective training is possible and produces measurable trade-offs.

Our work differs in direction. PrefGen designs objectives forward (given known objectives, train to optimize them). We propose inferring objectives backward (given standard methods, what implicit objectives do they optimize for). This distinction matters because backward inference enables post-hoc diagnosis of existing methods without multi-objective engineering.

\subsubsection{Benchmark Evaluation}

LiveCodeBench addresses temporal contamination through continuously updated problems. NaturalCodeBench reveals that 80\% of HumanEval tasks have low real-world correlation. These establish that existing metrics are imperfect. Our signature detection framework could provide diagnostic value for validating what benchmarks measure if the methodology validates with real models.
